\documentclass[../main.tex]{subfiles}

\begin{document}

\section{Task 2 --- Full Launch Vehicle Model}\label{sec:task2}

Task~2 reintroduces the full dynamics of Eq.~\eqref{eq:plant}: the lightly damped
first bending mode, the second-order TVC actuator of Eq.~\eqref{eq:tvc} with its
$20$~ms delay, and the INS measurement model of Eq.~\eqref{eq:ins}. I build the
design up incrementally, choosing the bending filter first and then --- unlike a
naive reuse of the Task-1 gains --- \textbf{re-tuning the PD on the full loop}:
the Task-1 PD was placed on the \emph{ideal}-actuator loop, and once the real TVC
servo, transport delay and notch are in the path their phase lag collapses the
rigid phase margin, so the gains must be re-tuned to recover it (D'Antuono: the
TVC actuator dynamics must be taken into account when placing the
PD~\cite{barrowsorr}).

\subsection{Step B --- uncompensated full model}

Adding the TVC, delay and bending mode to the Task~1 controller leaves the
low-frequency margins almost untouched, but the bending resonance shows up in
the loop at $|L(\omega_{BM})| = +29$~dB. With a damping of only
$\zeta_{BM} = 0.005$ this peak crosses the critical point and the closed loop
goes \textbf{unstable} (a closed-loop pole at $\mathrm{Re}(s)=+1.6$).

\subsection{Step C --- bending filter trade study}

The guideline filter of Eq.~4 admits two readings: as printed (negative
numerator damping, a non-minimum-phase \textbf{lead-lag} that phase-stabilises the
resonance by rotating it on the Nichols chart) or in its minimum-phase variant
(a band-reject \textbf{notch} that gain-stabilises the resonance by driving it
below $0$~dB)~\cite[Lecture~18]{dclvnotes}. Instead of picking one up front, I
compare four candidates at the fixed Task-1 gains (all built by
\texttt{build\_notch\_filter.m}, Listing~\ref{lst:hm3_notch});
Table~\ref{tab:filters} collects the resulting margins, classified by frequency
band (aerodynamic gain margin, rigid-body phase and gain margin, delay margin on
top of the $20$~ms already modelled, and the residual loop gain at the bending
frequency). The trade fixes the filter \emph{architecture}; the PD is re-tuned
afterwards.

\begin{itemize}
\item \textbf{C-LL --- Eq.~4 lead-lag alone.} Swept over the ranges suggested in
the assignment ($\zeta_{x,N} \in [0.1,0.3]$, $\zeta_{x,D} \in [0.4,0.6]$,
$\omega_x = \omega_{BM} \pm 4$~rad/s, $75$ combinations), $45/75$ now stabilise
the loop --- the milder $+29$~dB resonance (a consequence of the smaller Task-1
derivative gain) is within reach of pure phase rotation. But the best
combination ($\omega_x = 18.9$, $\zeta_{x,N}=0.30$, $\zeta_{x,D}=0.60$) is only
barely stable ($\mathrm{Re}(s) = -0.06$) with a thin rigid phase margin
($11^\circ$): phase stabilisation of the first mode leaves almost no room.
\item \textbf{C-N --- deep minimum-phase notch.} A deep min-phase notch
($\zeta_{x,N}=0.002$, $\zeta_{x,D}=0.7$) gain-stabilises the resonance, driving
$|L(\omega_{BM})|$ to $-22$~dB. The depth follows the closed-form attenuation of a
second-order notch, $K_\mathrm{notch}=20\log_{10}(\zeta_{x,N}/\zeta_{x,D})
=20\log_{10}(0.002/0.7)\approx-51$~dB~\cite[Lecture~18]{dclvnotes}, which is why
the sizing has to leave the assignment's suggested $\zeta_{x,N}\in[0.1,0.3]$
range. I keep the null narrow because the bending frequency sits only a decade
above the rigid crossover: a wide band-reject section would leak phase lag into
the conditionally-stable low-frequency region. This is the retained architecture.
\item \textbf{C-T --- notch triplet.} The notch-triplet recipe against
bending-frequency uncertainty places three cascaded notches at
$\{0.9,\,1.0,\,1.1\}\,\omega_{BM}$~\cite[Lecture~18]{dclvnotes}. Sized to the
same depth, the three wide ($\zeta_{x,D}=0.7$) sections pile up about $30^\circ$
of phase lag at the rigid crossover and the loop goes \textbf{unstable} at
nominal: the first mode is too close to the control band to blanket-notch
(triplets belong to the higher, well-separated modes).
\item \textbf{C-NLL --- notch $+$ lead-lag.} Reading the guideline literally
(``the Notch filter and possibly other filters''), the deep notch is cascaded
with the Eq.~4 lead-lag, the partner swept over the same ranges: $67/75$
combinations stabilise the loop. The best one ($\omega_x = 22.9$,
$\zeta_{x,N}=0.10$, $\zeta_{x,D}=0.40$, selected for maximum delay margin) is
stable but with clearly thinner nominal margins ($PM = 8.0^\circ$). Its payoff
is the bending-frequency robustness shown below.
\end{itemize}

\begin{table}[H]
\centering
\caption{Task 2 --- bending-filter trade at the fixed Task-1 gains (before the
full-loop re-tune). Margins classified by frequency band.}
\label{tab:filters}
\footnotesize
\setlength{\tabcolsep}{4pt}
\begin{tabular}{lcccccc}
\toprule
Candidate & Aero $|GM|$ [dB] & Rigid $PM$ [$^\circ$] & Rigid $|GM|$ [dB] & $|L(\omega_{BM})|$ [dB] & DM [ms] & stable \\
\midrule
B (no filter)      & 6.25 & 26.1  & ---   & $+29.0$ & ---  & no \\
C-LL (lead-lag)    & 5.96 & 11.4  & 8.02  & $+23.0$ & 15   & yes \\
C-N (deep notch)   & 6.08 & 14.6  & 10.43 & $-21.9$ & 98   & \textbf{yes} \\
C-T (triplet)      & ---  & $-7.3$ & ---  & $-55.8$ & ---  & no \\
C-NLL (notch+LL)   & 5.72 & 8.0   & 7.41  & $-28.1$ & 54   & yes \\
\bottomrule
\end{tabular}
\end{table}

\subsection{Retained design and PD re-tuning}

I keep the \textbf{deep notch}: it gain-stabilises the resonance
($|L(\omega_{BM})| = -22$~dB) and preserves the low-frequency aerodynamic gain
margin. But with the Task-1 gains --- placed on the \emph{ideal}-actuator loop
--- the combined phase lag of the second-order TVC, the $20$~ms transport delay
and the notch has eaten most of the rigid phase margin: it collapses from
$30^\circ$ to $\mathbf{14.6^\circ}$, with the delay margin down at the
$\sim100$~ms guideline (Table~\ref{tab:filters}, C-N). The design must therefore
be \textbf{re-tuned on the full loop}: raising the derivative gain restores the
phase lost to the lag. Solving for the $6$~dB / $30^\circ$ target on the full
open loop gives
\begin{equation}
K_{P,\theta} = 1.73,\qquad K_{D,\theta} = 0.69
\end{equation}
(the derivative gain up from the Task-1 $0.44$), which recovers a rigid phase
margin of $30^\circ$ at $3.2$~rad/s and a delay margin of $165$~ms, with the
bending still gain-stabilised at $|L(\omega_{BM})| = -18$~dB --- an $18$~dB
bending gain margin, above the $\ge 12$~dB usually required of gain-stabilised
modes~\cite{greensite_cr826}. Figure~\ref{fig:t2retune} shows the recovery on the
Nichols chart, and Figure~\ref{fig:t2nichols} the retained full-loop margins.

\begin{figure}[H]
\centering
\includegraphics[width=0.66\textwidth]{task2_retune.png}
\caption{Task 2 --- PD re-tuning on the full loop. With the Task-1 gains (ideal
actuator) the actuator $+$ delay $+$ notch lag drives the loop close to the
$+180^\circ$ critical point (rigid $PM = 15^\circ$, blue); re-tuning on the full
loop pulls the curve away and restores $PM = 30^\circ$ (orange).}
\label{fig:t2retune}
\end{figure}

\begin{figure}[H]
\centering
\includegraphics[width=0.64\textwidth]{task2_nichols.png}
\caption{Task 2 --- retained full-loop Nichols (deep notch, re-tuned PD) in the
launch-vehicle convention (critical point $+180^\circ$): aerodynamic gain margin
($|GM_a|=6$~dB), rigid phase margin ($30^\circ$) and the high-frequency rigid
gain margin (set by the actuator lag) are marked. The bending mode is
gain-stabilised, so its lobe stays well below $0$~dB near $\omega_{BM}$ rather
than threading the critical point.}
\label{fig:t2nichols}
\end{figure}

The filter trade itself is shown on Figure~\ref{fig:t2trade}: without
compensation the $+29$~dB resonance destabilises the loop (red); the best Eq.-4
lead-lag alone (green) now barely stabilises the milder peak but with a thin
margin; the deep notch drives the bending lobe far below $0$~dB (blue).

\begin{figure}[H]
\centering
\includegraphics[width=0.64\textwidth]{task2_nichols_trade.png}
\caption{Task 2 --- bending-filter trade on the Nichols chart (launch-vehicle
convention, critical point $+180^\circ$; a common phase shift aligns the three
loops at the rigid crossover): without a filter the $+29$~dB bending lobe (red)
wraps over the critical point and the loop is unstable; the best Eq.-4 lead-lag
alone (green) rotates the lobe but only marginally; the deep notch (blue) pushes
the bending lobe far below $0$~dB (the small loop near $-18$~dB), gain-stabilising
it.}
\label{fig:t2trade}
\end{figure}

The $s$-plane view of Figure~\ref{fig:t2pz} makes the mechanism explicit: without
the notch the closed-loop bending poles are pushed into the right half-plane by
the $+29$~dB resonance and the loop is unstable, whereas the deep notch drops the
loop gain at $\omega_{BM}$ and returns them to their lightly damped, stable
open-loop location just left of the imaginary axis. The notch gain-stabilises the
mode; it does not add damping to it.

\begin{figure}[H]
\centering
\includegraphics[width=0.7\textwidth]{task2_notch_poles.png}
\caption{Task 2 --- closed-loop bending poles without the notch (crosses, in the
right half-plane) and with the notch (circles, returned to the stable
neighbourhood of $\omega_{BM}$).}
\label{fig:t2pz}
\end{figure}

The root locus of the full loop (Figure~\ref{fig:t2rl}), from the built-in
\texttt{rlocus}, tells the same story as a gain sweep. Figure~\ref{fig:t2rl}(a)
zooms on the low-frequency rigid and lateral-drift dynamics; near $\omega_{BM}$
the bending branch skirts the imaginary axis while the notch keeps the
closed-loop poles ($k=1$) just inside the left half-plane (Figure~\ref{fig:t2pz});
and its far branches race toward the right-half-plane zeros of the Pad\'e delay
approximation, resolved in the full-scale view of Figure~\ref{fig:t2rl}(b).

\begin{figure}[H]
\centering
\begin{subfigure}{0.66\textwidth}
\centering
\includegraphics[width=\textwidth]{task2_rootlocus.png}
\caption{Zoom near the origin.}
\label{fig:t2rl_zoom}
\end{subfigure}

\vspace{0.6em}

\begin{subfigure}{0.66\textwidth}
\centering
\includegraphics[width=\textwidth]{task2_rootlocus_full.png}
\caption{Full view, including the Pad\'e-delay RHP zeros.}
\label{fig:t2rl_full}
\end{subfigure}
\caption{Task 2 --- root locus of the full loop (flexible plant, TVC servo,
Pad\'e delay and notch; $\times$: open-loop poles, $\circ$: open-loop zeros). The
design is stable, with all closed-loop poles in the left half-plane.
(a)~Low-frequency detail near the origin: the rigid-body and lateral-drift
dynamics, the unstable aerodynamic pole at $+1.8$~s$^{-1}$ and the real zeros at
$-2.5$ and $-15$~s$^{-1}$; the bending poles ($\omega_{BM}\approx 18.9$~rad/s),
the notch and the far delay zeros fall outside this window.
(b)~The same locus zoomed out: the far branches terminate on the transport-delay
RHP zeros at $+232$ and $+184\pm175i$~s$^{-1}$ (mirror LHP poles at $-232$ and
$-184\pm175i$).}
\label{fig:t2rl}
\end{figure}

\subsection{Sensitivity to the bending-frequency knowledge}

The deep notch pays for its performance with a $-51$~dB null only $0.08$~rad/s
wide, i.e.\ it assumes $\omega_{BM}$ is known exactly --- true in this exercise,
where the LPV dataset provides it, but never true in practice.
Table~\ref{tab:wbm} makes the trade explicit: each filter is held fixed (with the
re-tuned PD) while the true plant bending frequency is perturbed. The deep notch
tolerates $-10\%$ but goes unstable at $+5\%$ --- the asymmetric edge of its
narrow null; the notch$+$lead-lag combination, by phase-stabilising whatever
leaks past the detuned notch, holds the full $\pm 10\%$ band at the cost of the
thin nominal margins seen above. In a flight design, where $\omega_{BM}$ drifts
with propellant depletion, the robust variant (or a phase-stabilised loop) would
be mandatory; for this max-$\bar q$ snapshot I keep the exact notch.

\begin{table}[H]
\centering
\caption{Task 2 --- closed-loop stability with the filters fixed (re-tuned PD)
and the true bending frequency perturbed by up to $\pm 10\%$.}
\label{tab:wbm}
\begin{tabular}{lccccc}
\toprule
Candidate & $0.90\,\omega_{BM}$ & $0.95\,\omega_{BM}$ & $\omega_{BM}$ & $1.05\,\omega_{BM}$ & $1.10\,\omega_{BM}$ \\
\midrule
C-N --- deep notch        & yes & yes & yes & no  & no \\
C-NLL --- notch+lead-lag  & yes & yes & yes & yes & yes \\
\bottomrule
\end{tabular}
\end{table}

\subsection{Validation and comparison with the rigid model}

The gust response of the stabilised full model is shown in
Figure~\ref{fig:t2gust}; all four assignment time histories remain well damped,
with peaks $\theta = 0.23^\circ$, $z = 2.27$~m, $\delta = 0.51^\circ$ and a
load indicator $\bar q\,\alpha = 20.5$~kPa\,deg, essentially the rigid
case. Because the notch is localised at the bending frequency and the TVC
bandwidth ($70$~rad/s) is well above the control crossover, the full-model
response is almost indistinguishable from the rigid one
(Figure~\ref{fig:t2cmp}): the added dynamics affect stability robustness (the
bending and delay margins) far more than nominal performance.

\begin{figure}[H]
\centering
\includegraphics[width=0.85\textwidth]{task2_gust_response.png}
\caption{Task 2 --- full-model response to the severe wind gust: pitch attitude
$\theta$, lateral drift $z$, drift rate $\dot z$ and TVC deflection $\delta$.}
\label{fig:t2gust}
\end{figure}

\begin{figure}[H]
\centering
\includegraphics[width=0.95\textwidth]{task2_comparison_rigid_vs_full.png}
\caption{Task 2 --- rigid (Task 1) versus full-model gust response: pitch
attitude, lateral drift and TVC deflection are nearly coincident.}
\label{fig:t2cmp}
\end{figure}

\subsection{A note on the tail-wags-dog effect}

A second flexible effect bounds what any filter can achieve. The gimballed
nozzle has its own inertia, so its angular acceleration reacts back on the
airframe. Following the course treatment~\cite[Lecture~20]{dclvnotes}, the
oscillating nozzle exerts on the vehicle a lateral reaction force and a moment
\begin{equation}
R = -\,m_E \ell_E\,\ddot\delta, \qquad
M_{TWD} = \left(I_E + m_E \ell_E^{2}\right)\ddot\delta,
\label{eq:twd}
\end{equation}
proportional to the gimbal acceleration $\ddot\delta$ ($m_E$ nozzle mass,
$\ell_E$ pivot-to-nozzle-CG distance, $I_E$ nozzle inertia). This
\textbf{tail-wags-dog} (TWD) coupling feeds the second derivative of the control
back into the plant; it adds phase lag and an effective transmission zero that
caps the achievable closed-loop bandwidth. The
$-\phi_{1,TVC}T_c$ entry in the bending-forcing column of Eq.~\eqref{eq:plant} is
the flexible-mode analogue of this nozzle-reaction term. My model
treats the actuator as a stiff second-order block (Eq.~\eqref{eq:tvc}) and so
neglects the nozzle back-reaction --- the dog-wags-tail feedback, which
would add the gimbal angle as an extra state. A higher-fidelity design would
keep it, one of the extensions noted in the conclusions.

\end{document}
